\section{Domain dependence and polylogarithmic arithmetic cutoffs}
\label{p1:consequences:section}

Two features of the arithmetic estimates have uses beyond the dyadic
law theorem. Changing the interval of starting points changes the dominant
rational channel. Revealing more windows changes the set of available
prime coordinates, and the resulting loss must be measured by rank,
not just by the number of exceptional positions. We record these facts,
then explain how to cut off large ternary supports without assuming
uniformity of a fixed-order counting theorem.

\subsection{The aligned profile on a macroscopic interval}

For an interval of starts $A$, let $\sigma(x,y)$ be the dimension
of its rational even-unit code and put
\[
 R_{\rm rat}(A,B)=\sum_{\substack{x,y\in A\\|x-y|>L}}
             (2^{\sigma(x,y)}-1).
\]
The rational code is zero unless there are at least two exact units.
Let $R_{\rm al}(A,B)$ and $C_{\rm al}(A,L)$ be the sums over
aligned ordered separated pairs of $2^{\rho(x,y)}-1$ and
$\eta(x,y)2^{\rho(x,y)}-1$, respectively. Alignment has the
definition in Section~\ref{p1:arith:section}; $B=L+1$.

\begin{proposition}[Profiles depend on the interval]
\label{p1:consequences:macro}
In the dyadic critical band, writing $B=3k+r$, $0\le r<3$,
\begin{equation}
 R_{\rm rat}(I_N,B)\sim R_{\rm al}(I_N,B)
       \sim \frac{3+r}{18}NB2^k,
 \qquad
 R_{\rm al}(I_N,B)-R_{\rm rat}(I_N,B)
       =O_{C,\varepsilon}(N^{5/4+\varepsilon}).
 \label{p1:consequences:dyadic-positive}
\end{equation}
Fix $0<\delta<1$, let $A_{M,\delta}=[\lceil M^\delta\rceil,M)
\cap\mathbb Z$, and suppose $|L-\log_2M|\le C$.
For $B=2k+r$, $0\le r<2$,
\begin{equation}
 R_{\rm rat}(A_{M,\delta},B)\sim R_{\rm al}(A_{M,\delta},B)
       \sim\frac{2+r}{4}MB2^k,
 \label{p1:consequences:macro-positive}
\end{equation}
and their difference is $O_{C,\delta,\varepsilon}(M^{4/3+\varepsilon})$.
For the signed profile, write instead $L=2m+s$, $0\le s<2$.
Then
\begin{align}
 C_{\rm al}(A_{M,\delta},L)
   &=(M-2M^\delta)S_{1,2}(L)
           +O_{C,\delta,\varepsilon}(M^{4/3+\varepsilon}),
          \label{p1:consequences:macro-signed}\\
 S_{1,2}(L)
   &=2^{m-1}\big[(s+1)m+s^2+s+4\big]-3L+1
       \qquad(m\ge1),\label{p1:consequences:S12}\\
 C_{\rm al}(A_{M,\delta},L)
   &\sim\frac{s+1}{4}ML2^m.
          \label{p1:consequences:macro-equivalent}
\end{align}
\end{proposition}

\begin{proof}
In the channel $1:2$, each orientation permits
$M/2-M^\delta+O(B)$ translations of the unit parameter.
For $h=y-2x\in\{-1,\ldots,B-2\}$, the exact units have
the size of the congruence class of $h$ modulo two in this
interval. If these two sizes are $m_0,m_1$, their plateau
contribution is
\[
 \sum_{j=0}^1m_j(2^{m_j-1}-1)
       =2^{k-1}\big[(2+r)k+2r\big]-B.
\]
Outside the plateau, the unit count drops linearly with the
distance, so the sum of the weights is $O(2^{B/2})$.
Multiplying by both orientations proves the main term in
\eqref{p1:consequences:macro-positive}. The height-one
channel has no exact unit for separated windows; its residual
cost is $M^{1+o(1)}$. All other reduced heights are at least
three and contribute at most $M B^{O(1)}2^{L/3+o(L)}$.

At a fixed height, the full kernel differs from the rational
code only when the normalized windows contain a defective
integer at a threshold proportional to $B$. There are
$M^{1/2+o(1)}$ such hosts, by the square--smooth decomposition
in Lemma~\ref{p1:arith:defects}. At height two their weight
is $2^{L/2+o(L)}$, so the residual costs $M^{1+o(1)}$.
The proof of Lemma~\ref{p1:arith:alignedrank} applies uniformly
down to $x,y\ge M^\delta$: its Runge distance becomes
$\gg_\delta\log M/\log B$, while its redundancy remains
$O(B/\log B)$. Thus the other aligned heights have the
stated $M^{4/3+\varepsilon}$ total bound, including all
residual relations. This proves the comparison of the two
positive sums.

Retaining the predecessor sign replaces each placement factor
$L+1-D$ by $g_L(D)=L-1-D$. Hence the signed contribution
per translation is $S_{1,2}(L)$; the identity
\eqref{p1:arith:channelidentity} gives
\eqref{p1:consequences:S12}. The endpoints cost
$O(B^3 2^{L/2})$ and the other channels and residual kernels
have the same bounds in absolute value as above. This proves
\eqref{p1:consequences:macro-signed} and its equivalent.

For the dyadic assertion, the height-two channel is confined
to the endpoints. The channel $2:3$ has translation count
$N/3+O(B)$ over the two orientations, and three residue
classes of total size $B=3k+r$. Its plateau weight is
$2^{k-1}[(3+r)k+2r]-B$. The height-one residual costs
$N^{1+o(1)}$; the remaining heights are at least four and
cost $N^{5/4+o(1)}$. The fixed-height residual costs
$N^{5/6+o(1)}$. These estimates prove
\eqref{p1:consequences:dyadic-positive}.
\end{proof}

The residue classes of $B$ in the positive formulas and of
$L$ in the signed formulas are different. They cannot be
interchanged. With $\lambda_M=M2^{-L}$, the last formula gives
\[
 p^2C_{\rm al}(A_{M,\delta},L)
   \sim\frac{s+1}{4\,2^{s/2}}\lambda_M^{3/2}L M^{-1/2}.
\]
This is an aligned-pair statement. It does not assert a
macroscopic analogue of Theorem~\ref{p1:thm:main}. The actual
mean parameter for this interval would be $|A_{M,\delta}|p$,
whereas $\lambda_M$ expresses the ambient critical scale.

\subsection{How exceptional vertices bound the conditional rank}

Fix $j$ old windows and let $W$ be the square-class span of
all their raw values. Let $\mathcal P$ contain the primes at
most $B$ and every prime larger than $B$ dividing an old value.
Put
\[
 \mathcal S(\mathcal P)=\{n\le3N:\sf(n)\text{ is supported on }
                    \mathcal P\},\qquad
 t_x=|V_x\cap\mathcal S(\mathcal P)|.
\]

\begin{lemma}[Conditional rank loss]
\label{p1:consequences:conditional-rank}
Let $\rho_x$ be the dimension of the kernel of the map from
$E_x$ to square classes modulo $W$. Then
\[
 \rho_x\le(t_x-1)_+.
\]
Conditional on any realizable vector of the old raw signs,
the probability of a new start is either zero or $p2^{\rho_x}$.
It equals $p$ if $t_x\le1$.

Let $O$ be the set of starts at distance at most $L$ from an
old start, and put
$D_T=\#\{x\in I_N\setminus O:t_x>T\}$ for $1\le T\le L$.
Uniformly over the old signs,
\begin{equation}
 \E[Z_N\mid\mathcal F_{\rm old}]
   \le\lambda_N+|O|+p2^{T-1}B|\mathcal S(\mathcal P)|+D_T.
 \label{p1:consequences:conditional-budget}
\end{equation}
For fixed $j$ and $T=\lfloor(1-\delta)L\rfloor$, $0<\delta<1$,
the third term is $N^{1/2-\delta+o_{C,j}(1)}$.
\end{lemma}

\begin{proof}
A nonzero element of $E_x$ is an even nonempty selection of
values. If its square class lies in $W$, every prime larger
than $B$ with odd valuation in a selected value must divide
an old value: it cannot divide a second value of $V_x$.
Thus every selected value lies in $\mathcal S(\mathcal P)$.
The kernel injects into the even-subset space on these $t_x$
values, of dimension $(t_x-1)_+$. The conditional distribution
of the new relative signs is uniform on an affine image of
dimension $L-\rho_x$, proving the probability statement.

Use probability at most one on $O$ and on $\{t_x>T\}$.
On $\{t_x\le1\}$ the probability is $p$; on
$\{2\le t_x\le T\}$ it is at most $p2^{T-1}$.
Moreover $\#\{x:t_x\ge2\}\le\sum_xt_x
\le B|\mathcal S(\mathcal P)|$.
This proves \eqref{p1:consequences:conditional-budget}.
The Euler-product bound in Lemma~\ref{p1:limits:branching}
gives $|\mathcal S(\mathcal P)|=N^{1/2+o_{C,j}(1)}$ when
$j$ is fixed; substituting $T$ proves the final assertion.
\end{proof}

For example, if $Q(\mathcal P)$ counts pairs $n<m$ in
$\mathcal S(\mathcal P)$ with $m-n\le L$, double counting gives
\[
 D_T\le\frac{B Q(\mathcal P)}{\binom{T+1}{2}}.
\]
No independence of the two exceptional integers is assumed.
The rank factor explains why merely counting dependent windows
throws away useful conditional information.

\subsection{Cutting off large primes and large support groups}

Use the triples of blocks and the groups $G_u$ from
Proposition~\ref{p1:arith:ternary-budget}. The offsets and
width constants are fixed; a polynomial number in $B$ of
such patterns is allowed. Let $\mathcal Q$ be the set of
primes greater than the diameter threshold that occur with
odd valuation in all three blocks. A positive mass below
means a sum of nonnegative host weights, each at most
$N^{-\beta+o(1)}$, with $\beta$ fixed before $N$ and uniform
$o(1)$. For an original
primitive triple one can take $\beta=2$: two of its
window spaces are in direct sum.

\begin{lemma}[A Rankin bound for offset patterns]
\label{p1:consequences:rankin-patterns}
For fixed $c,D_0>0$, $d\in\{2,3\}$ and $X\ge1$, put
$s=1+(d-1)/c$. Then
\begin{equation}
 \sum_{\substack{Q\le X\ {\rm squarefree}\\P^-(Q)>B^c}}
       (D_0B)^{d\omega(Q)}
 \le X^s\exp\{O_{c,D_0}(B/\log B)\}.
 \label{p1:consequences:rankin-primes}
\end{equation}
If instead the factors are arbitrary distinct integers
$g>B^c$, then, for every fixed $0<\zeta<1$ and
$s_\zeta=1+(d-1+\zeta)/c$,
\begin{equation}
 \sum_{\substack{A\subset\mathbb N\ {\rm finite}\\
                g>B^c\ (g\in A),\ \prod_{g\in A}g\le X}}
       (D_0B)^{d|A|}
 \le X^{s_\zeta}\exp\{O_{c,d,\zeta,D_0}(B^{1-\zeta})\},
 \label{p1:consequences:rankin-groups}
\end{equation}
where $D_0$ is fixed.
\end{lemma}

\begin{proof}
Rankin's inequality bounds the first sum by
$X^s\prod_{q>B^c}(1+(D_0B)^dq^{-s})$.
Chebyshev's prime estimate and partial summation give
$\sum_{q>Y}q^{-s}\ll_s Y^{1-s}/\log Y$.
Since $d+c(1-s)=1$, taking the logarithm proves
\eqref{p1:consequences:rankin-primes}. For integer factors
the analogous product is
$\prod_{g>B^c}(1+(D_0B)^dg^{-s_\zeta})$.
Bound its logarithm by the corresponding sum over integers;
the exponent of $B$ is now $1-\zeta$.
\end{proof}

\begin{proposition}[Polylogarithmic confinement]
\label{p1:consequences:confinement}
For every fixed $c>0$ and $\varepsilon>0$, there is a fixed
integer $k$ such that, outside a mass
\begin{equation}
 O_{c,\beta,\varepsilon}
 \left(N^{-\min\{\beta-6/5,\,\beta-1-2/c\}+\varepsilon}\right),
 \label{p1:consequences:confinement-error}
\end{equation}
one has
\begin{equation}
 \#\{u:G_u>B^c\}
       \le k-1+\left\lfloor\frac{\log N}{c\log B}\right\rfloor=o(L).
 \label{p1:consequences:confinement-count}
\end{equation}
The same assertion holds for $\#\{q\in\mathcal Q:q>B^c\}$.
On a family of $B^{O(1)}$ fixed affine planes for the centers,
replace $2/c$ in the exceptional mass by $1/c$.
For original triples, $c>6$ is therefore sufficient for an
exceptional mass $o(LN^{-2/3})$; on these planes $c>3$
suffices. These are sufficient thresholds, not optimality claims.
\end{proposition}

\begin{proof}
First, the fixed-$k$ budget of
Proposition~\ref{p1:arith:ternary-budget} gives a mass
$O(N^{6/5-\beta+\varepsilon/3})$ for $\mathcal G_k>N$.
The identical budget holds for the product of the $k$ largest
individual primes. Indeed CRT fixes each center modulo their
product, exactly as in \eqref{p1:arith:groupcrt}. In the
three-prime tail, repeated offset supports cannot cause a
new case: if all supports coincide the product divides one
value; if exactly two coincide, or three are collinear, it
divides a nonzero coordinate of $d\times(X+u)$, of size
$O(NB)$. A zero cross product would give a forbidden pair
relation, as in the proof of the group budget. The large
product therefore forces a genuine plane. The CRT tail and
the induction on fixed $k$ apply unchanged.

Outside this first exception at most $k-1$ factors exceed
$N^{1/k}$. If the product of the remaining factors in
$(B^c,N^{1/k}]$ exceeds $N$, multiply them in a deterministic
order until their product first exceeds $N$. The resulting
$Q$ satisfies $N<Q\le X=N^{1+1/k}$.
For individual primes, choosing their three offsets costs
$(D_0B)^{3\omega(Q)}$ for a fixed width constant $D_0$.
CRT gives $O(1)$ choices for each center.
Equation~\eqref{p1:consequences:rankin-primes} bounds the
weighted mass by
\[
 N^{-\beta+(1+1/k)(1+2/c)+o(1)}.
\]
For groups, each factor has one support, costing at most
$(D_0B)^3$. Their pairwise coprimality is used in CRT first;
it may then be dropped in the positive sum over factor sets.
Equation~\eqref{p1:consequences:rankin-groups} gives instead
\[
 N^{-\beta+(1+1/k)(1+(2+\zeta)/c)+o(1)}.
\]
Choose $\zeta$ sufficiently small and then $k$ sufficiently
large, both before $N$, to absorb the extra exponent into
$\varepsilon$. If the centers lie on one of the specified
planes, enumerate only the two free coordinates and their
offsets; the third is fixed by the plane. This uses $d=2$
and gives $1/c$ in place of $2/c$. In the remaining
configurations the product is at most $N$, so the number
of factors greater than $B^c$ is at most
$\lfloor\log N/(c\log B)\rfloor$.
\end{proof}

The cutoff applies to whole support groups, even if a group
contains many primes. It can therefore be imposed before a
coordinate regularization. Reversing that order could merge
support labels and invalidate a bound stated only for the
original labels.
